Una vez entrenado el modelo, se evalúa su rendimiento sobre el subconjunto de evaluación llamando a la función \texttt{predict}. La evaluación del clasificador se valora mediante 3 métricas: la precisión, la matriz de confusión y el reporte de clasificación.

La precisión proporciona una visión sencilla del rendimiento del clasificador midiendo la proporción de datos correctamente clasificados con respecto al total de casos. No obstante, en problemas con un conjunto de datos desbalanceado como es este, puede reflejar un resultado engañoso, ya que se podría obtener un valor elevado debido a que predice correctamente las clases mayoritarias mientras que no lo hace para las minoritarias. Es por esto que se complementa esta métrica con las otras dos. Esta se obtiene mediante la función \texttt{accuracy\_score}

La matriz de confusión muestra como se reparten los aciertos y los fallos para cada clase. Las filas representan las clases reales y las columnas las clases predichas por el modelo, por lo que una celda cualquiera $(i,j)$ muestra cuántas clases cuya clase real es $i$ han sido predichas como $j$. Por esto mismo, la diagonal principal de la matriz mostrará los aciertos, mientras que cualquier valor fuera de esa diagonal es un error en la predicción. Esta es especialmente útil para analizar qué clases son confundidas entre sí. Esta se obtiene mediante la función \texttt{confusion\_matrix}.

El reporte de clasificación aporta para cada clase la precisión, exhaustividad y la medida F1. La precisión indica la proporción de las asignaciones del modelo a esa clase fue correctamente clasificada, mientras que la exhaustividad indica la proporción de las muestras de esa clase fue correctamente asignada por el clasificador. La medida F1 realiza la media armónica de estas 2, siendo la medida más real de precisión para valorar el rendimiento del modelo, especialmente en conjuntos de datos desbalanceados como este. Este se obtiene mediante la función \texttt{classification\_report}. Además, al final del reporte de clasificación se incluyen la media ponderada o \texttt{weighted avg} y la media macro o \texttt{macro avg}. La media ponderada calcula la media de las métricas ponderando por el número de muestras de cada clase, es decir, las clases mayoritarias tendrán más influencia en este resultado. Por otro lado, la media macro calcula la media de las métricas dando el mismo peso a todas las clases, es decir, sin tener en cuenta el número de muestras, reflejando el resultado con más precisión si el conjunto de datos es desbalanceado.

Finalmente, también se analizó la importancia de las características empleadas por el modelo. Consiste en un análisis que identifica qué características influyen en mayor medida en la toma de decisiones del clasificador, lo cual es de gran utilidad debido a que la primera fase de este trabajo fue definir unas características adicionales a las que ya contenía el conjunto de datos y se podrá valorar la incorporación de estas.